% Options for packages loaded elsewhere
\PassOptionsToPackage{unicode}{hyperref}
\PassOptionsToPackage{hyphens}{url}
%
\documentclass[
  ignorenonframetext,
  aspectratio=169,
  russian,
]{beamer}
\usepackage{pgfpages}
\setbeamertemplate{caption}[numbered]
\setbeamertemplate{caption label separator}{: }
\setbeamercolor{caption name}{fg=normal text.fg}
\beamertemplatenavigationsymbolshorizontal
% Prevent slide breaks in the middle of a paragraph
\widowpenalties 1 10000
\raggedbottom
\setbeamertemplate{part page}{
  \centering
  \begin{beamercolorbox}[sep=16pt,center]{part title}
    \usebeamerfont{part title}\insertpart\par
  \end{beamercolorbox}
}
\setbeamertemplate{section page}{
  \centering
  \begin{beamercolorbox}[sep=12pt,center]{section title}
    \usebeamerfont{section title}\insertsection\par
  \end{beamercolorbox}
}
\setbeamertemplate{subsection page}{
  \centering
  \begin{beamercolorbox}[sep=8pt,center]{subsection title}
    \usebeamerfont{subsection title}\insertsubsection\par
  \end{beamercolorbox}
}
\AtBeginPart{
  \frame{\partpage}
}
\AtBeginSection{
  \ifbibliography
  \else
    \frame{\sectionpage}
  \fi
}
\AtBeginSubsection{
  \frame{\subsectionpage}
}

\usepackage{amsmath,amssymb}
\usepackage{iftex}
\ifPDFTeX
  \usepackage[T1]{fontenc}
  \usepackage[utf8]{inputenc}
  \usepackage{textcomp} % provide euro and other symbols
\else % if luatex or xetex
  \usepackage{unicode-math}
  \defaultfontfeatures{Scale=MatchLowercase}
  \defaultfontfeatures[\rmfamily]{Ligatures=TeX,Scale=1}
\fi
\usepackage{lmodern}
\ifPDFTeX\else  
    % xetex/luatex font selection
\fi
% Use upquote if available, for straight quotes in verbatim environments
\IfFileExists{upquote.sty}{\usepackage{upquote}}{}
\IfFileExists{microtype.sty}{% use microtype if available
  \usepackage[]{microtype}
  \UseMicrotypeSet[protrusion]{basicmath} % disable protrusion for tt fonts
}{}
\usepackage{xcolor}
\newif\ifbibliography
\setlength{\emergencystretch}{3em} % prevent overfull lines
\setcounter{secnumdepth}{-\maxdimen} % remove section numbering


\providecommand{\tightlist}{%
  \setlength{\itemsep}{0pt}\setlength{\parskip}{0pt}}\usepackage{longtable,booktabs,array}
\usepackage{calc} % for calculating minipage widths
\usepackage{caption}
% Make caption package work with longtable
\makeatletter
\def\fnum@table{\tablename~\thetable}
\makeatother
\usepackage{graphicx}
\makeatletter
\newsavebox\pandoc@box
\newcommand*\pandocbounded[1]{% scales image to fit in text height/width
  \sbox\pandoc@box{#1}%
  \Gscale@div\@tempa{\textheight}{\dimexpr\ht\pandoc@box+\dp\pandoc@box\relax}%
  \Gscale@div\@tempb{\linewidth}{\wd\pandoc@box}%
  \ifdim\@tempb\p@<\@tempa\p@\let\@tempa\@tempb\fi% select the smaller of both
  \ifdim\@tempa\p@<\p@\scalebox{\@tempa}{\usebox\pandoc@box}%
  \else\usebox{\pandoc@box}%
  \fi%
}
% Set default figure placement to htbp
\def\fps@figure{htbp}
\makeatother
% definitions for citeproc citations
\NewDocumentCommand\citeproctext{}{}
\NewDocumentCommand\citeproc{mm}{%
  \begingroup\def\citeproctext{#2}\cite{#1}\endgroup}
\makeatletter
 % allow citations to break across lines
 \let\@cite@ofmt\@firstofone
 % avoid brackets around text for \cite:
 \def\@biblabel#1{}
 \def\@cite#1#2{{#1\if@tempswa , #2\fi}}
\makeatother
\newlength{\cslhangindent}
\setlength{\cslhangindent}{1.5em}
\newlength{\csllabelwidth}
\setlength{\csllabelwidth}{3em}
\newenvironment{CSLReferences}[2] % #1 hanging-indent, #2 entry-spacing
 {\begin{list}{}{%
  \setlength{\itemindent}{0pt}
  \setlength{\leftmargin}{0pt}
  \setlength{\parsep}{0pt}
  % turn on hanging indent if param 1 is 1
  \ifodd #1
   \setlength{\leftmargin}{\cslhangindent}
   \setlength{\itemindent}{-1\cslhangindent}
  \fi
  % set entry spacing
  \setlength{\itemsep}{#2\baselineskip}}}
 {\end{list}}
\usepackage{calc}
\newcommand{\CSLBlock}[1]{\hfill\break\parbox[t]{\linewidth}{\strut\ignorespaces#1\strut}}
\newcommand{\CSLLeftMargin}[1]{\parbox[t]{\csllabelwidth}{\strut#1\strut}}
\newcommand{\CSLRightInline}[1]{\parbox[t]{\linewidth - \csllabelwidth}{\strut#1\strut}}
\newcommand{\CSLIndent}[1]{\hspace{\cslhangindent}#1}

\IfFileExists{plex-otf.sty}{
  %% Full TeXlive
  % \usepackage[%
  %   % math,
  %   RM={Scale=0.94},SS={Scale=0.94},SScon={Scale=0.94},TT={Scale=MatchLowercase,FakeStretch=0.9},DefaultFeatures={Ligatures=Common}
  % ]{plex-otf}
}{
  %% TinyTeX
  \usepackage{libertine}
}

%%% Load theme
% https://deic.uab.cat/~iblanes/beamer_gallery/
\IfFileExists{beamerthemegotham.sty}{
  %% Full TeXlive
  \usetheme{gotham}
  \gothamset{
    numbering=totalpagenumber,
    parttocframe default=off,
    sectiontocframe default=off,
    subsectiontocframe default=off,
  }
}{
  %% TinyTeX
  \usetheme{Madrid}
}
\makeatletter
\@ifpackageloaded{caption}{}{\usepackage{caption}}
\AtBeginDocument{%
\ifdefined\contentsname
  \renewcommand*\contentsname{Содержание}
\else
  \newcommand\contentsname{Содержание}
\fi
\ifdefined\listfigurename
  \renewcommand*\listfigurename{Список иллюстраций}
\else
  \newcommand\listfigurename{Список иллюстраций}
\fi
\ifdefined\listtablename
  \renewcommand*\listtablename{Список таблиц}
\else
  \newcommand\listtablename{Список таблиц}
\fi
\ifdefined\figurename
  \renewcommand*\figurename{Рисунок}
\else
  \newcommand\figurename{Рисунок}
\fi
\ifdefined\tablename
  \renewcommand*\tablename{Таблица}
\else
  \newcommand\tablename{Таблица}
\fi
}
\@ifpackageloaded{float}{}{\usepackage{float}}
\floatstyle{ruled}
\@ifundefined{c@chapter}{\newfloat{codelisting}{h}{lop}}{\newfloat{codelisting}{h}{lop}[chapter]}
\floatname{codelisting}{Список}
\newcommand*\listoflistings{\listof{codelisting}{Листинги}}
\makeatother
\makeatletter
\makeatother
\makeatletter
\@ifpackageloaded{caption}{}{\usepackage{caption}}
\@ifpackageloaded{subcaption}{}{\usepackage{subcaption}}
\makeatother

\ifLuaTeX
\usepackage[bidi=basic]{babel}
\else
\usepackage[bidi=default]{babel}
\fi
\babelprovide[main,import]{russian}
\babelprovide[import]{english}
% get rid of language-specific shorthands (see #6817):
\let\LanguageShortHands\languageshorthands
\def\languageshorthands#1{}
\usepackage{csquotes}
\usepackage{bookmark}

\IfFileExists{xurl.sty}{\usepackage{xurl}}{} % add URL line breaks if available
\urlstyle{same} % disable monospaced font for URLs
\hypersetup{
  pdftitle={Лабораторная работа № 5},
  pdfauthor={Гайдук Софья Сергеевна},
  pdflang={ru-RU},
  hidelinks,
  pdfcreator={LaTeX via pandoc}}


\title{Лабораторная работа № 5}
\subtitle{Отчет}
\author{Гайдук Софья Сергеевна}
\date{}

\begin{document}
\frame{\titlepage}

\renewcommand*\contentsname{Содержание}
\begin{frame}[allowframebreaks]
  \frametitle{Содержание}
  \setcounter{tocdepth}{1}
  \tableofcontents
\end{frame}

\section{Цель
работы}\label{ux446ux435ux43bux44c-ux440ux430ux431ux43eux442ux44b}

\begin{frame}{Цель работы}
Установка и настройка менеджеров паролей pass. Управление файлами
конфигурации.
\end{frame}

\section{Выполнение лабораторной
работы}\label{ux432ux44bux43fux43eux43bux43dux435ux43dux438ux435-ux43bux430ux431ux43eux440ux430ux442ux43eux440ux43dux43eux439-ux440ux430ux431ux43eux442ux44b}

\begin{frame}{Выполнение лабораторной работы}
Установка pass (рис.~\ref{fig-001}).

\begin{figure}

\centering{

\includegraphics[width=0.7\linewidth,height=\textheight,keepaspectratio]{image/image1.jpg}

}

\caption{\label{fig-001}image1}

\end{figure}%

Установка gopass (рис.~\ref{fig-002}).

\begin{figure}

\centering{

\includegraphics[width=0.7\linewidth,height=\textheight,keepaspectratio]{image/image2.jpg}

}

\caption{\label{fig-002}image2}

\end{figure}%

Просмотр списка ключей (рис.~\ref{fig-003}).

\begin{figure}

\centering{

\includegraphics[width=0.7\linewidth,height=\textheight,keepaspectratio]{image/image3.jpg}

}

\caption{\label{fig-003}image3}

\end{figure}%

Инициализация хранилища (рис.~\ref{fig-004}).

\begin{figure}

\centering{

\includegraphics[width=0.7\linewidth,height=\textheight,keepaspectratio]{image/image4.jpg}

}

\caption{\label{fig-004}image4}

\end{figure}%

Создадим структуру git (рис.~\ref{fig-005}).

\begin{figure}

\centering{

\includegraphics[width=0.7\linewidth,height=\textheight,keepaspectratio]{image/image5.jpg}

}

\caption{\label{fig-005}image5}

\end{figure}%

Создадим репозиторий (рис.~\ref{fig-006}).

\begin{figure}

\centering{

\includegraphics[width=0.7\linewidth,height=\textheight,keepaspectratio]{image/image6.jpg}

}

\caption{\label{fig-006}image6}

\end{figure}%

Зададим адрес репозитория на хостинге (рис.~\ref{fig-007}).

\begin{figure}

\centering{

\includegraphics[width=0.7\linewidth,height=\textheight,keepaspectratio]{image/image7.jpg}

}

\caption{\label{fig-007}image7}

\end{figure}%

Для синхронизации выполним следующие команды (рис.~\ref{fig-008}).

\begin{figure}

\centering{

\includegraphics[width=0.7\linewidth,height=\textheight,keepaspectratio]{image/image8.jpg}

}

\caption{\label{fig-008}image8}

\end{figure}%

Вручную закоммитим и выложим изменения (рис.~\ref{fig-009}).

\begin{figure}

\centering{

\includegraphics[width=0.7\linewidth,height=\textheight,keepaspectratio]{image/image9.jpg}

}

\caption{\label{fig-009}image9}

\end{figure}%

Проверим статус синхронизации модной командой (рис.~\ref{fig-010}).

\begin{figure}

\centering{

\includegraphics[width=0.7\linewidth,height=\textheight,keepaspectratio]{image/image10.jpg}

}

\caption{\label{fig-010}image10}

\end{figure}%

Установим программу, обеспечивающую интерфейс native messaging
(рис.~\ref{fig-011}).

\begin{figure}

\centering{

\includegraphics[width=0.7\linewidth,height=\textheight,keepaspectratio]{image/image11.jpg}

}

\caption{\label{fig-011}image11}

\end{figure}%

Установим интерфейс для взаимодействия с броузером (native
messaging)(рис.~\ref{fig-012}, рис.~\ref{fig-013}).

\begin{figure}

\centering{

\includegraphics[width=0.7\linewidth,height=\textheight,keepaspectratio]{image/image12.jpg}

}

\caption{\label{fig-012}image12}

\end{figure}%

\begin{figure}

\centering{

\includegraphics[width=0.7\linewidth,height=\textheight,keepaspectratio]{image/image13.jpg}

}

\caption{\label{fig-013}image13}

\end{figure}%

Добавим новый пароль (рис.~\ref{fig-014}).

\begin{figure}

\centering{

\includegraphics[width=0.7\linewidth,height=\textheight,keepaspectratio]{image/image14.jpg}

}

\caption{\label{fig-014}image14}

\end{figure}%

Отобразим пароль для указанного имени файла (рис.~\ref{fig-015}).

\begin{figure}

\centering{

\includegraphics[width=0.7\linewidth,height=\textheight,keepaspectratio]{image/image15.jpg}

}

\caption{\label{fig-015}image15}

\end{figure}%

Заменим существующий пароль (рис.~\ref{fig-016}).

\begin{figure}

\centering{

\includegraphics[width=0.7\linewidth,height=\textheight,keepaspectratio]{image/image16.jpg}

}

\caption{\label{fig-016}image16}

\end{figure}%

Установим дополнительное программное обеспечение (рис.~\ref{fig-017}).

\begin{figure}

\centering{

\includegraphics[width=0.7\linewidth,height=\textheight,keepaspectratio]{image/image17.jpg}

}

\caption{\label{fig-017}image17}

\end{figure}%

Установим шрифты (рис.~\ref{fig-018}, рис.~\ref{fig-019},
рис.~\ref{fig-020}).

\begin{figure}

\centering{

\includegraphics[width=0.7\linewidth,height=\textheight,keepaspectratio]{image/image18.jpg}

}

\caption{\label{fig-018}image18}

\end{figure}%

\begin{figure}

\centering{

\includegraphics[width=0.7\linewidth,height=\textheight,keepaspectratio]{image/image19.jpg}

}

\caption{\label{fig-019}image19}

\end{figure}%

\begin{figure}

\centering{

\includegraphics[width=0.7\linewidth,height=\textheight,keepaspectratio]{image/image20.jpg}

}

\caption{\label{fig-020}image20}

\end{figure}%

Установим бинарный файл (рис.~\ref{fig-021}).

\begin{figure}

\centering{

\includegraphics[width=0.7\linewidth,height=\textheight,keepaspectratio]{image/image21.jpg}

}

\caption{\label{fig-021}image21}

\end{figure}%

Создадим свой репозиторий для конфигурационных файлов на основе шаблона
(рис.~\ref{fig-022}).

\begin{figure}

\centering{

\includegraphics[width=0.7\linewidth,height=\textheight,keepaspectratio]{image/image22.jpg}

}

\caption{\label{fig-022}image22}

\end{figure}%

Инициализируем chezmoi с нашим репозиторием dotfiles
(рис.~\ref{fig-023}).

\begin{figure}

\centering{

\includegraphics[width=0.7\linewidth,height=\textheight,keepaspectratio]{image/image23.jpg}

}

\caption{\label{fig-023}image23}

\end{figure}%

Проверим, какие изменения внесёт chezmoi в домашний каталог
(рис.~\ref{fig-024}).

\begin{figure}

\centering{

\includegraphics[width=0.7\linewidth,height=\textheight,keepaspectratio]{image/image24.jpg}

}

\caption{\label{fig-024}image24}

\end{figure}%

Если устраивают изменения, внесённые chezmoi, запустим
(рис.~\ref{fig-025}).

\begin{figure}

\centering{

\includegraphics[width=0.7\linewidth,height=\textheight,keepaspectratio]{image/image25.jpg}

}

\caption{\label{fig-025}image25}

\end{figure}%

На второй машине инициализируем chezmoi с нашим репозиторием dotfiles
(рис.~\ref{fig-026}).

\begin{figure}

\centering{

\includegraphics[width=0.7\linewidth,height=\textheight,keepaspectratio]{image/image26.jpg}

}

\caption{\label{fig-026}image26}

\end{figure}%

Проверим, какие изменения внесёт chezmoi в домашний каталог
(рис.~\ref{fig-027}).

\begin{figure}

\centering{

\includegraphics[width=0.7\linewidth,height=\textheight,keepaspectratio]{image/image27.jpg}

}

\caption{\label{fig-027}image27}

\end{figure}%

Если устраивают изменения, внесённые chezmoi, запустим
(рис.~\ref{fig-028}).

\begin{figure}

\centering{

\includegraphics[width=0.7\linewidth,height=\textheight,keepaspectratio]{image/image28.jpg}

}

\caption{\label{fig-028}image28}

\end{figure}%

При существующем каталоге chezmoi получим и применим последние изменения
из нашего репозитория (рис.~\ref{fig-029}).

\begin{figure}

\centering{

\includegraphics[width=0.7\linewidth,height=\textheight,keepaspectratio]{image/image29.jpg}

}

\caption{\label{fig-029}image29}

\end{figure}%

Установим свои dotfiles на новый компьютер с помощью одной команды
(рис.~\ref{fig-030}).

\begin{figure}

\centering{

\includegraphics[width=0.7\linewidth,height=\textheight,keepaspectratio]{image/image30.jpg}

}

\caption{\label{fig-030}image30}

\end{figure}%

Извлекем последние изменения из репозитория и применим их
(рис.~\ref{fig-031}).

\begin{figure}

\centering{

\includegraphics[width=0.7\linewidth,height=\textheight,keepaspectratio]{image/image31.jpg}

}

\caption{\label{fig-031}image31}

\end{figure}%

Извлекем последние изменения из своего репозитория и посмотрим, что
изменится, фактически не применяя изменения (рис.~\ref{fig-032}).

\begin{figure}

\centering{

\includegraphics[width=0.7\linewidth,height=\textheight,keepaspectratio]{image/image32.jpg}

}

\caption{\label{fig-032}image32}

\end{figure}%

Мы довольны изменениями и можем применить их (рис.~\ref{fig-033}).

\begin{figure}

\centering{

\includegraphics[width=0.7\linewidth,height=\textheight,keepaspectratio]{image/image33.jpg}

}

\caption{\label{fig-033}image33}

\end{figure}%

Автоматически фиксируем и отправляем изменения в репозиторий
(рис.~\ref{fig-034}).

\begin{figure}

\centering{

\includegraphics[width=0.7\linewidth,height=\textheight,keepaspectratio]{image/image34.jpg}

}

\caption{\label{fig-034}image34}

\end{figure}%
\end{frame}

\section{Выводы}\label{ux432ux44bux432ux43eux434ux44b}

\begin{frame}{Выводы}
Мы научились установке и настройке менеджеров паролей pass, а также
управлению файлами конфигурации.
\end{frame}

\section*{Список
литературы}\label{ux441ux43fux438ux441ux43eux43a-ux43bux438ux442ux435ux440ux430ux442ux443ux440ux44b}
\addcontentsline{toc}{section}{Список литературы}

\begin{frame}{Список литературы}
1.Kulyabov. Лабораторная работа № 5. Настройка рабочей среды. RUDN

\phantomsection\label{refs}
\begin{CSLReferences}{0}{1}
\end{CSLReferences}
\end{frame}




\end{document}
